\chapter{Infertility}

\section*{Definition}
Infertility is defined as the inability to achieve pregnancy after 12 months of regular, unprotected intercourse in women under age 35, or after 6 months for women age 35 and older. Primary infertility refers to couples who have never conceived, while secondary infertility applies to those who have previously achieved a pregnancy. Recurrent pregnancy loss (two or more failed clinical pregnancies) is considered a distinct but related condition.

\section*{What Happens in Your Body}
Fertility requires a complex sequence of ovulation, gamete transport, fertilization, and implantation. Disruption at any step causes infertility: anovulatory disorders impair oocyte release (hypothalamic-pituitary-ovarian axis dysfunction), tubal pathology prevents sperm-egg meeting, endometriosis creates a hostile peritoneal environment, male factor abnormalities reduce sperm quantity or quality, and uterine factors interfere with implantation. Shared genetic and epigenetic mechanisms involving DNA repair genes and hormonal signaling pathways contribute across categories.

\section*{Causes}
\begin{itemize}
  \item Female factor (approximately 35\%): Ovulatory disorders (PCOS 80\% of anovulation, hypothalamic amenorrhea, hyperprolactinemia, primary ovarian insufficiency), Tubal disease (salpingitis, pelvic adhesions, prior ectopic, tubal ligation), Endometriosis, Uterine factors (fibroids, polyps, Asherman syndrome, congenital anomalies), Diminished ovarian reserve, Cervical factors (stenosis, hostile mucus)
  \item Male factor (approximately 30\%): Abnormal semen parameters (oligospermia, asthenospermia, teratospermia), Varicocele, Hormonal disorders (hypogonadotropic hypogonadism), Obstructive azoospermia (CFTR mutations, vasectomy), Genetic causes (Klinefelter syndrome, microdeletions of Y-chromosome), Antisperm antibodies, Idiopathic male infertility
  \item Combined factors (approximately 20\%): Both partners have contributing abnormalities
  \item Unexplained infertility (approximately 15\%): Standard testing fails to identify a cause; likely involves subtle gamete dysfunction or implantation failure
\end{itemize}

\section*{When to See a Doctor}
See your doctor if:
\begin{itemize}
  \item Failure to conceive after 12 months if woman $<$ 35 years, or after 6 months if $\ge$ 35 years
  \item Known risk factor: irregular or absent menses, known endometriosis, prior pelvic infection, prior chemotherapy or pelvic radiation
  \item Male with history of cryptorchidism, testicular trauma/surgery, chemotherapy, or known genetic abnormality
  \item More than one first-trimester pregnancy loss
  \item Either partner with known genetic disorder that could affect fertility
\end{itemize}

\section*{Related Symptoms}
\begin{itemize}
  \item Oligomenorrhea or amenorrhea
  \item Dysmenorrhea
  \item Dyspareunia
  \item Hirsutism
  \item Galactorrhea
  \item Erectile dysfunction
  \item Ejaculatory dysfunction
  \item Decreased libido
  \item Testicular pain or atrophy
  \item Recurrent pregnancy loss
\end{itemize}
\textbf{Important:} Both partners should be evaluated simultaneously with a coordinated approach including semen analysis, ovulation assessment, tubal patency testing, and ovarian reserve testing.

\section*{Self-Care / Home Management}
\begin{itemize}
  \item Optimize timing of intercourse using ovulation predictor kits (urinary LH surge) or fertility awareness methods targeting the fertile window (6-day period ending on ovulation day).
  \item Both partners should maintain healthy BMI (18.5--24.9), avoid tobacco, limit alcohol, and eliminate recreational drug use.
  \item Prenatal vitamin with at least 400 mcg folic acid daily for the female partner.
  \item Limit caffeine intake to less than 200 mg daily (approximately one 12-oz coffee).
  \item Avoid excessive exercise that may disrupt menstrual cyclicity.
\end{itemize}

\section*{How Your Doctor Will Evaluate You}
\begin{itemize}
  \item Female evaluation includes menstrual history, ovulation confirmation (mid-luteal progesterone, urine LH kits, or basal body temperature charting), transvaginal ultrasound, anti-M\"ullerian hormone, and day 3 FSH, LH, estradiol.
  \item Hysterosalpingography (HSG) assesses tubal patency and uterine cavity; saline infusion sonohysterography (SIS) evaluates intracavitary lesions.
  \item Male evaluation begins with semen analysis (two samples spaced 2--12 weeks apart) assessing concentration, motility, and morphology per WHO criteria.
  \item Additional testing: TSH, prolactin, hemoglobin A1c, karyotype (both partners), and cystic fibrosis carrier screening.
  \item Genetic carrier screening may identify actionable mutations before pursuing assisted reproductive technology.
\end{itemize}

\section*{Treatment Options}
\begin{itemize}
  \item Ovulation induction with clomiphene citrate or letrozole for anovulatory infertility; metformin added for PCOS with insulin resistance.
  \item Intrauterine insemination (IUI) with or without controlled ovarian hyperstimulation for unexplained, mild male factor, or cervical factor infertility.
  \item In vitro fertilization (IVF) with or without intracytoplasmic sperm injection (ICSI) for tubal factor, severe male factor, endometriosis, or failed IUI.
  \item Surgical correction: TESA/TESE/MESA for obstructive azoospermia, varicocelectomy for clinically significant varicocele, hysteroscopic resection of polyps/fibroids, laparoscopic treatment of endometriosis.
  \item Third-party reproduction (donor sperm, donor oocytes, gestational carrier) for cases where autologous gametes are not viable.
  \item Preimplantation genetic testing (PGT) for couples at risk of transmitting known genetic disorders.
\end{itemize}

\section*{References}

\begin{thebibliography}{99}
\bibitem{harrison-infert} Longo DL, et al. \emph{Harrison's Principles of Internal Medicine}. 21st ed. McGraw-Hill; 2022. Chapter 400: Infertility.
\bibitem{uptodate-infert} Practice Committee of the American Society for Reproductive Medicine. Infertility evaluation in women. In: UpToDate, Post TW (Ed), UpToDate, Waltham, MA; 2025.
\bibitem{asrm} Practice Committee of the ASRM. Definitions of infertility and recurrent pregnancy loss. \emph{Fertil Steril}. 2023;120(1):54-63.
\bibitem{nice-infert} National Institute for Health and Care Excellence. \emph{Fertility Problems: Assessment and Treatment}. NICE Guideline [CG156]; 2023.
\bibitem{speroff} Speroff L, et al. \emph{Clinical Gynecologic Endocrinology and Infertility}. 10th ed. Wolters Kluwer; 2023. Chapter 8: Infertility.
\bibitem{krausz} Krausz C, et al. Male infertility: a clinical review. \emph{Lancet}. 2023;401(10377):718-731.
\end{thebibliography}
